% Required packages:
% \usepackage{amsmath,amssymb}
% \usepackage{xcolor}

\subsection{Derivation of the Information--Response Bound}
\label{app:information_response_bound}

This appendix derives the bound
\begin{equation}
\boxed{
T_\pi(n)
\ge
\frac{2a_n(1-a_n)}{\ln 2}\,
\chi_{h_c}(n/N)^2
}
\label{eq:IR_bound_final}
\end{equation}
used to define the information--response efficiency.
The derivation proceeds in three steps:
(i) expressing the conditional mutual information as a weighted
Jensen--Shannon divergence,
(ii) applying Pinsker's inequality, and
(iii) bounding the change in the mean target fraction by total variation.

Throughout the derivation, the current population state \(n\) is held fixed.

% ------------------------------------------------------------
\subsubsection{Conditional mutual information as a Jensen--Shannon divergence}

Let the binary control variable be
\[
U\in\{0,1\},
\]
where \(U=0\) denotes no intervention and \(U=1\) denotes advocacy of
the controller target \(Z\). At population state \(n\), define
\begin{equation}
\Pr(U=1\mid n)=a_n,
\qquad
\Pr(U=0\mid n)=1-a_n.
\end{equation}

The corresponding next-round target-count distributions are
\begin{align}
R_0(n'\mid n)
&=
\Pr(n'\mid U=0,n),
\\
R_{h_c}(n'\mid n)
&=
\Pr(n'\mid U=1,n).
\end{align}

For notational convenience, while \(n\) is fixed, write
\begin{equation}
P(n')\equiv R_0(n'\mid n),
\qquad
Q(n')\equiv R_{h_c}(n'\mid n),
\qquad
a\equiv a_n.
\label{eq:PQ_definitions}
\end{equation}

The conditional mutual information between the control action and the
next population state is, in the round-indexed notation of the main text,
\[
T_\pi(n)\equiv I(U_k;n_{k+1}\mid n_k=n).
\]
Holding \(n_k=n\) fixed, we use the shorthand \(U\equiv U_k\) and
\(n'\equiv n_{k+1}\) below. Thus
\begin{align}
I(U;n'\mid n)
&=
\sum_{u,n'}
\Pr(u,n'\mid n)
\log_2
\frac{
\Pr(u,n'\mid n)
}{
\Pr(u\mid n)\Pr(n'\mid n)
}.
\label{eq:conditional_mi_start}
\end{align}

Using the factorization
\[
\Pr(u,n'\mid n)
=
\Pr(u\mid n)\Pr(n'\mid u,n),
\]
the factor \(\Pr(u\mid n)\) cancels inside the logarithm:
\begin{align}
I(U;n'\mid n)
&=
\sum_{u,n'}
\Pr(u\mid n)\Pr(n'\mid u,n)
\log_2
\frac{
\Pr(n'\mid u,n)
}{
\Pr(n'\mid n)
}
\nonumber\\
&=
\sum_u
\Pr(u\mid n)
D_{\mathrm{KL}}
\left(
\Pr(n'\mid u,n)
\,\Vert\,
\Pr(n'\mid n)
\right).
\label{eq:mi_as_average_kl}
\end{align}

We therefore need the next-state distribution when the value of \(U\)
is not specified. By the law of total probability,
\begin{align}
\Pr(n'\mid n)
&=
\sum_{u\in\{0,1\}}
\Pr(u\mid n)\Pr(n'\mid u,n)
\nonumber\\
&=
(1-a)P(n')+aQ(n').
\end{align}

Define the mixture distribution
\begin{equation}
M(n')
\equiv
(1-a)P(n')+aQ(n').
\label{eq:mixture_M}
\end{equation}

Substituting the two possible actions into
Eq.~\eqref{eq:mi_as_average_kl} gives
\begin{align}
I(U;n'\mid n)
&=
(1-a)D_{\mathrm{KL}}(P\Vert M)
+
aD_{\mathrm{KL}}(Q\Vert M).
\label{eq:mi_two_kl}
\end{align}

The weighted Jensen--Shannon divergence is defined by
\begin{equation}
\operatorname{JS}_{a}(P,Q)
\equiv
(1-a)D_{\mathrm{KL}}(P\Vert M)
+
aD_{\mathrm{KL}}(Q\Vert M),
\qquad
M=(1-a)P+aQ.
\end{equation}

Hence, exactly,
\begin{equation}
\boxed{
I(U;n'\mid n)
=
\operatorname{JS}_{a}(P,Q)
}
\label{eq:mi_equals_js_short}
\end{equation}
or, restoring the round-indexed notation,
\begin{equation}
\boxed{
T_\pi(n)
\equiv
I(U_k;n_{k+1}\mid n_k=n)
=
\operatorname{JS}_{a_n}
\left[
R_0(\cdot\mid n),
R_{h_c}(\cdot\mid n)
\right].
}
\label{eq:mi_equals_js}
\end{equation}

Up to this point, all relations are exact; no bound has yet been used.

% ------------------------------------------------------------
\subsubsection{Pinsker bound on the Jensen--Shannon divergence}

For two probability distributions \(A\) and \(B\), define their total
variation distance as
\begin{equation}
\operatorname{TV}(A,B)
\equiv
\frac{1}{2}
\sum_x
\left|A(x)-B(x)\right|.
\label{eq:tv_definition}
\end{equation}

Pinsker's inequality states that, when the Kullback--Leibler divergence
is measured in bits,
\begin{equation}
D_{\mathrm{KL}}(A\Vert B)
\ge
\frac{2}{\ln 2}\,
\operatorname{TV}(A,B)^2.
\label{eq:pinsker_bits}
\end{equation}
The factor \(1/\ln 2\) appears because the standard Pinsker inequality
is conventionally written for natural logarithms, whereas here the
information is reported in bits.

Applying Eq.~\eqref{eq:pinsker_bits} separately to the two KL terms in
Eq.~\eqref{eq:mi_two_kl} yields
\begin{align}
T_\pi(n)
&\ge
\frac{2}{\ln 2}
\left[
(1-a)\operatorname{TV}(P,M)^2
+
a\operatorname{TV}(Q,M)^2
\right].
\label{eq:pinsker_js_intermediate}
\end{align}

The two total variation distances simplify because \(M\) is a mixture
of \(P\) and \(Q\). From
\[
M-P
=
(1-a)P+aQ-P
=
a(Q-P),
\]
we obtain
\begin{align}
\operatorname{TV}(P,M)
&=
\frac12
\sum_{n'}
\left|
a\bigl(Q(n')-P(n')\bigr)
\right|
\nonumber\\
&=
a\,\operatorname{TV}(P,Q).
\label{eq:tv_PM}
\end{align}

Similarly,
\[
Q-M
=
Q-(1-a)P-aQ
=
(1-a)(Q-P),
\]
and therefore
\begin{equation}
\operatorname{TV}(Q,M)
=
(1-a)\operatorname{TV}(P,Q).
\label{eq:tv_QM}
\end{equation}

Substituting Eqs.~\eqref{eq:tv_PM}--\eqref{eq:tv_QM} into
Eq.~\eqref{eq:pinsker_js_intermediate},
\begin{align}
T_\pi(n)
&\ge
\frac{2}{\ln 2}
\left[
(1-a)a^2
+
a(1-a)^2
\right]
\operatorname{TV}(P,Q)^2
\nonumber\\
&=
\frac{2}{\ln 2}
a(1-a)
\underbrace{
\left[a+(1-a)\right]
}_{=1}
\operatorname{TV}(P,Q)^2.
\end{align}

Thus
\begin{equation}
\boxed{
T_\pi(n)
\ge
\frac{2a(1-a)}{\ln 2}
\operatorname{TV}(P,Q)^2.
}
\label{eq:js_tv_bound}
\end{equation}

Restoring \(a=a_n\), \(P=R_0\), and \(Q=R_{h_c}\),
\begin{equation}
T_\pi(n)
\ge
\frac{2a_n(1-a_n)}{\ln 2}
\operatorname{TV}
\left(
R_0(\cdot\mid n),
R_{h_c}(\cdot\mid n)
\right)^2.
\label{eq:js_tv_original}
\end{equation}

% ------------------------------------------------------------
\subsubsection{From total variation to the mean target-fraction response}

Let \(n'\in\{0,\ldots,N\}\) denote the next-round number of agents
supporting the controller target. Define the corresponding target
fraction
\begin{equation}
f(n')
\equiv
\frac{n'}{N}.
\label{eq:f_definition}
\end{equation}
Since \(0\le n'\le N\),
\begin{equation}
0\le f(n')\le 1.
\label{eq:f_bounded}
\end{equation}

Define the intervention-induced change in the mean target fraction by
\begin{align}
\chi_{h_c}(n/N)
&\equiv
\mathbb{E}_{R_{h_c}}
\left[
\frac{n'}{N}
\middle| n
\right]
-
\mathbb{E}_{R_0}
\left[
\frac{n'}{N}
\middle| n
\right]
\nonumber\\
&=
\mathbb{E}_{Q}[f]
-
\mathbb{E}_{P}[f].
\label{eq:chi_definition}
\end{align}

For completeness, we now show that any function satisfying
\(0\le f\le1\) obeys
\begin{equation}
\left|
\mathbb{E}_{Q}[f]
-
\mathbb{E}_{P}[f]
\right|
\le
\operatorname{TV}(P,Q).
\label{eq:bounded_expectation_tv}
\end{equation}

Define
\[
\Delta(n')\equiv Q(n')-P(n').
\]
Since both \(P\) and \(Q\) are normalized,
\begin{equation}
\sum_{n'}\Delta(n')=0.
\end{equation}

Let
\[
A_+
=
\left\{
n':\Delta(n')\ge0
\right\}.
\]
Then
\begin{equation}
\sum_{n'\in A_+}\Delta(n')
=
\operatorname{TV}(P,Q).
\label{eq:positive_tv}
\end{equation}

Using Eq.~\eqref{eq:chi_definition},
\begin{align}
\chi_{h_c}(n/N)
&=
\sum_{n'}
f(n')\Delta(n')
\nonumber\\
&=
\sum_{n'\in A_+}
f(n')\Delta(n')
+
\sum_{n'\notin A_+}
f(n')\Delta(n').
\end{align}

The second sum is non-positive because \(f(n')\ge0\) and
\(\Delta(n')<0\) outside \(A_+\). Therefore
\begin{align}
\chi_{h_c}(n/N)
&\le
\sum_{n'\in A_+}
f(n')\Delta(n')
\nonumber\\
&\le
\sum_{n'\in A_+}
\Delta(n')
\nonumber\\
&=
\operatorname{TV}(P,Q),
\end{align}
where the second inequality uses \(f(n')\le1\).
Repeating the argument after exchanging \(P\) and \(Q\) gives the
corresponding lower bound. Hence
\begin{equation}
\boxed{
|\chi_{h_c}(n/N)|
\le
\operatorname{TV}(P,Q).
}
\label{eq:chi_tv}
\end{equation}

Squaring both sides,
\begin{equation}
\chi_{h_c}(n/N)^2
\le
\operatorname{TV}(P,Q)^2.
\label{eq:chi_tv_squared}
\end{equation}

Finally, combining Eqs.~\eqref{eq:js_tv_bound} and
\eqref{eq:chi_tv_squared},
\begin{align}
T_\pi(n)
&\ge
\frac{2a_n(1-a_n)}{\ln 2}
\operatorname{TV}
\left(
R_0(\cdot\mid n),
R_{h_c}(\cdot\mid n)
\right)^2
\nonumber\\[1mm]
&\ge
\frac{2a_n(1-a_n)}{\ln 2}
\chi_{h_c}(n/N)^2.
\end{align}

Therefore,
\begin{equation}
\boxed{
T_\pi(n)
\ge
\frac{2a_n(1-a_n)}{\ln 2}\,
\chi_{h_c}(n/N)^2.
}
\end{equation}

For reference, the complete chain of relations is
\begin{equation}
\boxed{
\begin{aligned}
T_\pi(n)
&=
I(U_k;n_{k+1}\mid n_k=n)
\\
&=
\operatorname{JS}_{a_n}
\left[
R_0(\cdot\mid n),
R_{h_c}(\cdot\mid n)
\right]
\\
&\ge
\frac{2a_n(1-a_n)}{\ln 2}
\operatorname{TV}
\left(
R_0(\cdot\mid n),
R_{h_c}(\cdot\mid n)
\right)^2
\\
&\ge
\frac{2a_n(1-a_n)}{\ln 2}
\chi_{h_c}(n/N)^2.
\end{aligned}
}
\label{eq:complete_IR_chain}
\end{equation}
